\begin{textsample}{NEG Dim 3 – vem\_ed\_08 – Score 1.00 – vem\_ed\_08/t033.txt}  \label{ex:f3_neg_087}
Olhar atento ao tempo Maira de Lourdes Rezende, doutora em Nanociências pela UFABC (2017), é pesquisadora e professora no Curso Superior de Tecnologia em Polímeros da Fatec Sorocaba.
No segundo semestre de 2020, participou da oficina Taller COIL Diseño, ofertada por um consórcio de universidades chilenas do qual a Universidad Tecnológica de Chile (INACAP) faz parte.
A partir de uma aproximação proposta pelas coordenações das duas instituições, Maira e Patricio Hernán Marabolí Albornoz, professor do curso de Engenharia de Minas da INACAP, desenvolveram na oficina o PCI"Soluções alternativas para a reciclagem dos resíduos gerados na indústria extrativa mineira a partir de um trabalho colaborativo e intercultural".
Desse PCI, participaram 12 estudantes brasileiros e 12 chilenos.
A seguir, o depoimento da professora sobre essa experiência: Elaboramos um projeto interdisciplinar (Engenharia de Minas e Tecnologia em Polímeros) em consonância com os Objetivos de Desenvolvimento Sustentável (ODS) da Organização das Nações Unidas (ONU).
Dessa forma, definimos a reciclagem de resíduos da indústria extrativa mineira como foco do projeto.
O planejamento se baseou no cronograma de aulas das duas IES, para desenvolver semanalmente as atividades (síncronas e assíncronas).
Vídeos, artigos científicos e distintas metodologias de ensino contribuíram para atingir o principal objetivo do PCI: proporcionar vivência intercultural e promover desenvolvimento profissional e pessoal dos estudantes.
O"quebra-gelo"(rompehielo, em espanhol) ocorreu nas duas primeiras semanas: na primeira, os alunos deveriam escolher 5 fotografias (suas ou não) que os representassem.
Nessa ocasião, os grupos se reuniram para conhecer os colegas e criar vínculos.
Na segunda atividade, aplicamos um questionário para identificar a"orientação cultural"dos estudantes, que compartilharam suas respostas entre os membros do grupo e elaboraram um vídeo em que cada integrante abordou um tópico do questionário.
A participação foi ativa, excelente.
O principal desafio, logo no início do PCI, foi a comunicação entre os alunos e a organização do tempo para entregar as tarefas.
Ferramentas como o WhatsApp facilitam a interação.
Além disso, realizar reuniões periódicas para esclarecer dúvidas, bem como o alinhamento e o engajamento dos docentes, são extremamente eficientes para incentivar o cumprimento das tarefas.
A avaliação ocorreu em todas as etapas, com elaboração de textos, apresentações ou vídeos pelas equipes.
Ao fim do PCI, os grupos relataram em vídeo a experiência, suas dificuldades e êxitos.
Para realizar um PCI deve-se considerar idioma da instituição parceira, objetivos de aprendizagem e detalhar cronograma e metodologias.
Aconselho àqueles professores que se interessem em realizar um PCI, que o façam.
A experiência adquirida, tanto no âmbito pessoal quanto no profissional, é incrível.

% matched lemmas: 
\end{textsample}
